% page 12 of the facsimile (image p-011 of the scan)
% block 160.0 x 233.3 mm ; left margin 8.0 mm
\pgc{-12.700mm}{0.000mm}{
\l{\cel{5}4}
\l{}
\l{\sou{abrotano}.\cel{2}(\sou{bot}.) Speco di artemizio di qua la folii,}
\l{\cel{3}kande onu manuagas li, lasas sur la fingri odoro ci-}
\l{\cel{3}tronala. - L.\cel{1}\sou{artemisia abrotanum}.}
\l{}
\l{\sou{abrupta}. I. Di qua la pento esas ruptita bruske, ne-egale,}
\l{\cel{3}- II. (\sou{metaf.}) (\sou{pri stilo}).Di qua la aluro esas ne-}
\l{\cel{3}harmonioza e ne-egala. - DEFIS.}
\l{}
\l{\sou{absenta}. Qua ne esas en la loko ube lu povas o devas esar.}
\l{\cel{3}- DEFIRS.}
\l{}
\l{\sou{absinto}. (\sou{bot.})\cel{2}Planto aromata tre bitra, ek la familio}
\l{\cel{3}\textquotedbl{}kompozaji\textquotedbl{}. L.\cel{1}\sou{artemisia absinthium}. - DEFIRS.}
\l{}
\l{\sou{absoluta}.\cel{2}Qua ne submisesas a kondiciono, a restrikto. DEFIRS.}
\l{}
\l{\sou{absolvar}. (\sou{trans.}) Dispensar de la puniso, de la responso}
\l{\cel{3}a qua lu su\cel{1}\sou{exp}ozis pro ago pri qua lego preskriptas}
\l{\cel{3}puniso. - DEFIS.}
\l{}
\l{\sou{absorbar}.\cel{2}(\sou{trans}.) Igar penetrar, retenar en sua substanco,}
\l{\cel{3}liquido, gaso. - DEFIS.}
\l{}
\l{\sou{abstenar}. (\sou{trans.}) I. Interdiktar a su ipsa agar (ulo).-}
\l{\cel{3}II. Interdiktar a su ipsa uzar, juar (ulo). - DEFI.}
\l{}
\l{\sou{abstersar}. (\sou{trans.}) (\sou{medicino})\cel{2}Netigar de la puso, de la}
\l{\cel{3}sanio qua adheras a la surfaco di plago, di ulcero.DEFIS.}
\l{}
\l{\sou{abstinencar}. (\sou{netrans}.) (\sou{religio})\cel{2}Ne manjar karno en la dii}
\l{\cel{3}quin indikas la Eklezio katolika. - DEFIS.}
\l{}
\l{\sou{abstraktar}. (\sou{trans}.) I. Izolar per penso. - II. (\sou{filoz}.)}
\l{\cel{3}Konsiderar parto, izolante lu ek la toto. -\cel{1}\sou{Nombro abs}-}
\l{\cel{3}\sou{traktita}\cel{1}: qua ne indikas grandoro reala, objekti mezur-}
\l{\cel{3}ebla o kontebla. - DEFIRS.}
\l{}
\l{\sou{abstruza}.\cel{2}Desfacile komprenebla e tedanta. - DEFIS.}
\l{}
\l{\sou{absurda}.\cel{2}Kontrea ye la raciono komuna. - DEFIRS.}
\l{}
\l{\sou{abulio}.\cel{2}(\sou{medicino}). Morbo mentala quan karakterizas la}
\l{\cel{3}savo di to quon onu devas agar e la nepovo agar lo.DEFIS.}
\l{}
\l{\sou{abundar}.\cel{2}(\sou{netrans.})\cel{2}Esar plu multa kam suficante. - EFIS.}
\l{}
\l{\sou{abutar}. (\sou{netrans}.)Arivar per sua extremajo. (Ex.: strado}
\l{\cel{3}qua abutas an placo ; nervi qui abutas an la cerebro).EFIS.}
\l{}
\l{\sou{abutmento}. (\sou{teknol.})\cel{2}Konstrukturo, masivo petra, destinita}
\l{\cel{3}a suportar pulso. (Ex.: sur la rivi di fluvio, por rezis-}
\l{\cel{3}tar la pulsi da la arki). - F.}
}
